\chapter{Introducción}
\label{sec:intro}


\section{Contexto y Motivaciones}

\newpage
%\section{Alcance}

\section{Objetivos}

El objetivo principal de este TFG es el diseño y desarrollo de
un videojuego completamente funcional utilizando el motor de desarrollo Godot y el lenguaje de
programación C\#. El proyecto tiene como finalidad implementar un sistema completo
de juego que incluya las mecánicas básicas, un oponente controlado por inteligencia
artificial y una interfaz de usuario que permita al jugador interactuar de forma
clara e intuitiva con el sistema.

Para alcanzar este objetivo general, se han definido los siguientes objetivos
específicos:

\begin{itemize}
	
	\item Diseñar e implementar la lógica fundamental del videojuego, incluyendo la
	estructura del tablero, la definición de las piezas y las reglas de movimiento
	y combate entre ellas.
	
	\item Desarrollar un sistema de control de turnos y gestión del estado del juego
	que permita ejecutar partidas completas respetando las reglas definidas.
	
	\item Implementar un bot inteligente capaz de
	analizar el estado del tablero y seleccionar movimientos válidos durante la
	partida.
	
	\item Diseñar e implementar una interfaz de usuario clara e intuitiva que
	proporcione al jugador información visual sobre el estado del juego y permita
	interactuar con el sistema de manera sencilla.
	
	\item Incorporar elementos audiovisuales, como animaciones y efectos sonoros,
	que mejoren la experiencia de juego y aporten mayor claridad a las acciones
	realizadas durante la partida.
	
	\item Adquirir experiencia práctica en el uso del motor de videojuegos Godot y
	en el desarrollo de sistemas interactivos utilizando el lenguaje de programación
	C\#.
	
\end{itemize}

\section{Estructura del documento}
Este documento se organiza en distintos capítulos, que dividen el proyecto en ... 
